\section{Uniform Gaussian survival}\label{sec:survival}

We now prove Proposition~\ref{prop:filtered}.  The matrix estimates of the
preceding section measure how much the filtered constraints change as the
sequence index varies.  At a fixed spectral scale, we use that estimate
to control each bounded-energy group of states over one short interval
of Ornstein--Uhlenbeck time with a fixed positive probability.
Gaussian correlation then combines these intervals and the
spectral scales.  Their total cost per unit time is bounded by
$C\tr Q$, which gives the required Dirichlet-energy bound.

\subsection{Stationary survival and Dirichlet energy}

Let $Q$ be a positive definite diagonal matrix on $\R^m$, and let $U(\tau)$
be the stationary solution of
\[
 dU=-QU\,d\tau+\sqrt{2Q}\,dB(\tau).
\]
Here $B$ is standard $m$-dimensional Brownian motion and $U(0)$ has standard
Gaussian law $\gamma_m$, independently of its future driving noise.  Thus
the generator is
\[
 \mathcal L_Q=\tr(Q\nabla^2)-\langle Qx,\nabla\rangle,
 \qquad
 \operatorname{Cov}(U(\tau),U(\sigma))=e^{-Q|\tau-\sigma|}.
\]
Throughout, $\tau$ denotes Ornstein--Uhlenbeck time and $t$ denotes the
sequence index.

\begin{lemma}\label{lem:ou-energy}
Let $D\subset\R^m$ be bounded, nonempty, open and convex.  Define
\[
 \lambda_Q^{\mathrm{OU}}(D)
 =\inf_{0\ne h\in H_0^1(D)}
 \frac{\int_D\langle Q\nabla h,\nabla h\rangle\,d\gamma_m}
      {\int_D h^2\,d\gamma_m}.
\]
For every $T>0$,
\begin{equation}\label{eq:ou-comparison}
 \mathbb P\{U([0,T])\subset D\}
 \le e^{-T\lambda_Q^{\mathrm{OU}}(D)},
 \qquad
 \lambda_Q(D)\le\lambda_Q^{\mathrm{OU}}(D)+\tfrac12\tr Q.
\end{equation}
\end{lemma}

\begin{proof}
The Gaussian density and its reciprocal are bounded on $D$, so the
Gaussian and Lebesgue $H_0^1$ form domains agree.  Integration by parts
identifies the displayed form with the killed Ornstein--Uhlenbeck
semigroup $P_T^D$.  Its $L^2(\gamma_m)$ operator norm is
$e^{-T\lambda_Q^{\mathrm{OU}}(D)}$.  The killed-semigroup representation
used in Lemma~\ref{lem:brownian-rate} therefore gives
\[
 \mathbb P\{U([0,T])\subset D\}
 =\langle\mathbf1_D,P_T^D\mathbf1_D\rangle_{L^2(\gamma_m)}
 \le\gamma_m(D)e^{-T\lambda_Q^{\mathrm{OU}}(D)}.
\]
This proves the first inequality.  Equivalently, the representation follows
from the Dirichlet Feynman--Kac formula after the conjugation below; its
potential is bounded on $D$ \cite{DaviesSimon1984}.

Write $\rho(x)=(2\pi)^{-m/2}e^{-|x|^2/2}$ and put $g=h\sqrt\rho$.
For $h\in C_c^\infty(D)$, expansion and integration by parts give
\[
 \int_D\langle Q\nabla h,\nabla h\rangle\,d\gamma_m
 =\int_D\left[
    \langle Q\nabla g,\nabla g\rangle
    +\left(\tfrac14 x^TQx-\tfrac12\tr Q\right)g^2
   \right]dx.
\]
The identity extends by form closure to $H_0^1(D)$.  Multiplication by
$\sqrt\rho$ maps this space bijectively to itself and preserves the
corresponding $L^2$ norms.  Since the potential is at least
$-\tr Q/2$, taking Rayleigh infima proves the second inequality.
\end{proof}

\subsection{One energy interval at one spectral scale}

Fix the data of Proposition~\ref{prop:filtered} and a positive diagonal
$Q$.  We use the matrices and scalar quantities from the preceding
section:
\[
 L_t=R_tQR_t^T,\qquad X_t=L_t/\theta,\qquad
 F_t=F_t^\theta=R_t^Tf_0(X_t)R_t,\qquad
 k_t=k_t(\theta)=\tr\phi(X_t).
\]
In this subsection the dyadic scale $\theta$ is fixed.  Recall that
$a=3/2$ and that $E_\theta(s,t)$ is additive and nonnegative for $s\le t$.
Put $h_t=E_\theta(0,t)$.

\begin{lemma}\label{lem:epoch-prob}
There is an absolute $c_1$ with the following property.  Let
$\mathcal I\subset\{0,\ldots,n\}$ be an index set whose values $h_t$
lie in an interval of length at most $e_0$, where $e_0$ is the constant
in Lemma~\ref{lem:increment}.  For every $u\ge0$,
\begin{equation}\label{eq:epoch-prob}
 \mathbb P\left\{
   \lVert F_tU(\tau)\rVert\le c_1\sqrt{k_t}
   \text{ for every }t\in\mathcal I,
   \ u\le\tau\le u+\theta^{-1}
 \right\}\ge\tfrac12.
\end{equation}
The same conclusion holds on any shorter time interval.
\end{lemma}

\begin{proof}
We may omit indices with $F_t=0$.  By
\eqref{eq:cutoff-moments}, every remaining index has $k_t\ge1/4$.
Divide these indices into the classes
\[
 \mathcal I_j=\{t\in\mathcal I:K_j\le k_t<2K_j\},
 \qquad K_j=2^j/4,\qquad j=0,1,\ldots.
\]
We first bound the expected supremum within one nonempty class, writing
$k=K_j$.

Stationarity, coisometry and \eqref{eq:cutoff-moments} give
\begin{equation}\label{eq:vector-second-moment}
 \mathbb E\lVert F_tU(\tau)\rVert^2
 =\tr f_0(X_t)^2\le4k_t\le8k.
\end{equation}
For two indices in the class, Lemma~\ref{lem:increment}, with the indices
put in chronological order, implies
\begin{equation}\label{eq:index-increment}
 \mathbb E\lVert(F_t-F_s)U(\tau)\rVert^2
 \le Ck|h_t-h_s|^{1/a}.
\end{equation}
Here we dropped the positive right weight in that lemma and used
$1+k_t\le6k$.  For time increments at a fixed index,
\[
 \operatorname{Cov}(U(\tau)-U(\sigma))
 =2(I-e^{-Q|\tau-\sigma|})\preceq2|\tau-\sigma|Q,
\]
so that
\begin{align}
 \mathbb E\lVert F_t(U(\tau)-U(\sigma))\rVert^2
 &\le2|\tau-\sigma|\tr(F_t^2Q)\notag\\
 &=2\theta|\tau-\sigma|\tr(X_tf_0(X_t)^2)
 \le Ck\theta|\tau-\sigma|.\label{eq:time-increment}
\end{align}
This trace identity uses $R_tR_t^T=I$; it does not require $Q$ and
$F_t$ to commute.

We give the elementary chaining argument for the two parameters.  Choose
$H$ so that $h_t\in[H,H+e_0]$ for every index under consideration, and put
$\alpha=(h_t-H)/e_0$ and $\beta=\theta(\tau-u)$ in $[0,1]$.
Split a mixed increment into an
index increment and a time increment.  Equations~\eqref{eq:index-increment}
and \eqref{eq:time-increment} bound its second moment by
\[
 Ck\bigl(|\alpha-\alpha'|^{1/a}+|\beta-\beta'|\bigr),
\]
where $\alpha$ is the normalized energy coordinate; the constant absorbs
$e_0\le1$.  At depth $l$, divide both parameter intervals into dyadic
cells of width $2^{-l}$ and choose a representative from each occupied
product cell.  There are at most $C4^l$ representatives.  The increment
from each representative to that of its parent cell is a centered
Gaussian vector with second moment at most
\[
 v_l=Ck2^{-l/a}.
\]

For use here, if $G_1,\ldots,G_M$ are arbitrarily dependent centered
Gaussian vectors and $\mathbb E\lVert G_i\rVert^2\le v$, then
\begin{equation}\label{eq:gaussian-vector-max}
 \mathbb E\max_{i\le M}\lVert G_i\rVert
 \le2\sqrt{v(\log M+1/2)}.
\end{equation}
Indeed, for $v>0$, diagonalizing the covariance of $G_i$ gives
$\mathbb E e^{\lVert G_i\rVert^2/(4v)}\le e^{1/2}$, because its
eigenvalues are nonnegative and sum to at most $v$.  Bound the exponential
of the maximum by the sum of the exponentials, then use Jensen's
inequality and Cauchy--Schwarz.  The case $v=0$ is immediate.

Apply \eqref{eq:gaussian-vector-max} to the increments at each depth.
Starting with \eqref{eq:vector-second-moment} and summing the parent
increments gives an absolute $C_0$ such that
\begin{equation}\label{eq:class-mean}
 \mathbb E\sup_{\substack{t\in\mathcal I_j\\
                         u\le\tau\le u+\theta^{-1}}}
          \lVert F_tU(\tau)\rVert
 \le C\sqrt k\left(1+\sum_{l\ge1}
          2^{-l/(2a)}(1+\sqrt l)\right)
 \le C_0\sqrt k.
\end{equation}
To justify convergence of the representatives, note that there are only
finitely many distinct energy values.  If $h_t=h_s$,
Lemma~\ref{lem:increment} gives $F_t=F_s$ exactly.  At sufficiently fine
depth a target index therefore has the same filter as its representative,
and the representative times converge to its target time.  Continuity
of $U$ completes the passage to the supremum.  Equivalently, one may
first use finite time grids and then increase them to a dense set.
Neither argument uses a lower bound on the positive gaps between energy
values.  In particular, the number of indices does not enter
\eqref{eq:class-mean}.

It remains to control all the $k$ classes at once.  On a finite time
grid the supremum $S_j$ in \eqref{eq:class-mean} is a maximum of Gaussian
vector norms.  Represent the joint Gaussian vectors as linear images of
one standard Gaussian vector.  Each corresponding linear map has norm at
most one because its output covariance is $F_t^2\preceq I$, by
\eqref{eq:cutoff-moments}; consequently
$S_j$ is a $1$-Lipschitz function of that standard Gaussian vector.
Equivalently, it is a supremum of scalar Gaussian functionals, each
with variance at most one.

Markov's inequality and \eqref{eq:class-mean} give
$\mathbb P\{S_j\le2C_0\sqrt{K_j}\}\ge1/2$.
Borell's Gaussian enlargement inequality
\cite[Theorem~3.1]{Borell1975}, applied to this sublevel set, gives
\[
 \mathbb P\{S_j>2C_0\sqrt{K_j}+r\}
 \le1-\Phi(r)\le e^{-r^2/2},\qquad r\ge0,
\]
where $\Phi$ is the standard normal distribution function.  The same
bound holds for the full time supremum by increasing dense grids and
path continuity.  Choose $c_1>2C_0$.  The probability that any constraint in class
$j$ exceeds $c_1\sqrt{k_t}$ is at most
\[
 \exp\bigl(-(c_1-2C_0)^2K_j/2\bigr).
\]
With $A=(c_1-2C_0)^2/8$, the sum of these bounds is at most
\[
 \sum_{j\ge0}e^{-A2^j}
 \le\frac{e^{-A}}{1-e^{-A}}\le\tfrac12
\]
once $c_1$ is an absolute constant large enough that $A\ge\log3$.
This proves \eqref{eq:epoch-prob}.  Restricting the time interval can
only increase its probability.
\end{proof}

The two bounds in this proof have different roles: the trace of the
covariance controls the expected vector norm, while its operator norm
gives the absolute concentration parameter.  This distinction is what
allows the class failure probabilities to be summed without a factor
depending on the dimension.

\subsection{Combining the constraints}

\begin{proof}[Proof of Proposition~\ref{prop:filtered}]
Fix $d,n,Q$ and the entire recursion.  Let $\Theta$ be the finite set of
dyadic scales for which some $F_t^\theta$ is nonzero.  If $q_{\min}$ and
$q_{\max}$ are the smallest and largest diagonal entries of $Q$, then
coisometry gives $q_{\min}I\preceq L_t\preceq q_{\max}I$, and the
support of $f_0$ gives
\begin{equation}\label{eq:used-scales}
 \Theta\subset[q_{\min}/4,2q_{\max}],
 \qquad \sum_{\theta\in\Theta}\theta\le4q_{\max}\le4\tr Q.
\end{equation}
At each scale, partition the sequence indices by the half-open intervals
\[
 [r e_0,(r+1)e_0),\qquad r=0,1,\ldots,
\]
containing $E_\theta(0,t)$.  These are the energy intervals to which
Lemma~\ref{lem:epoch-prob} applies.  There are at most
$C(1+E_\theta(0,n))$ nonempty intervals, including any repetitions of
an energy value in the same interval.

For each fixed $t$, the scale constraints control the original state
norm.  To see this, diagonalize $L_t$ and write $y=R_tU(\tau)$ in this
eigenbasis, with eigenvalues $\lambda_i>0$.  Each $\lambda_i$ admits a
dyadic scale with $\lambda_i/\theta\in[1,2]$, where $f_0=1$; that scale
belongs to $\Theta$.  Since $R_t^T$ is an isometry,
\begin{align}
 \sum_{\theta\in\Theta}\lVert F_t^\theta U(\tau)\rVert^2
 &=\sum_i |y_i|^2\sum_{\theta\in\Theta}
                   f_0(\lambda_i/\theta)^2
 \ge\lVert R_tU(\tau)\rVert^2.\label{eq:spectral-reconstruction}
\end{align}
Thus all the scale constraints at a given time imply, by
\eqref{eq:scale-size},
\begin{equation}\label{eq:state-survival-bound}
 \lVert R_tU(\tau)\rVert^2
 \le c_1^2\sum_{\theta\in\Theta}k_t(\theta)
 \le c_1^2 C_\phi d
 \qquad(0\le t\le n),
\end{equation}
with an absolute $C_\phi$.

We also impose coordinate bounds.  A stationary scalar
Ornstein--Uhlenbeck process of rate $q>0$, restricted to an interval of
length at most $1/q$, stays in $[-2,2]$ with probability at least an
absolute $p_0>0$.  By stationarity and time rescaling, it suffices to
consider rate one on $[0,1]$.  In distribution that process is
\[
 e^{-s}\bigl(Z+B(e^{2s}-1)\bigr),\qquad 0\le s\le1,
\]
where $Z$ is standard normal and $B$ is an independent standard Brownian
motion.  The event $|Z|\le1$ and
$\sup_{r\le e^2-1}|B(r)|\le1$ has a fixed positive probability and
implies the asserted bound.

Choose absolute constants
\[
 A_0>c_1\sqrt{C_\phi},\qquad B_0>2,
\]
and let $\mathcal D$ be the domain in Proposition~\ref{prop:filtered}
with these constants.  It is bounded, open, centrally symmetric and
convex, and contains a neighborhood of zero.  The closed constraints
\eqref{eq:state-survival-bound} and $|U_j(\tau)|\le2$ lie strictly inside
$\mathcal D$.

For a fixed $T>0$, cover $[0,T]$ by at most $1+T\theta$ consecutive
closed intervals of length at most $1/\theta$ at scale $\theta$.
For every such interval and every nonempty energy interval impose the
event in Lemma~\ref{lem:epoch-prob}.  Also cover $[0,T]$ by at most
$1+TQ_{jj}$ intervals of length at most $1/Q_{jj}$ for coordinate $j$
and impose its box event.  All these events are closed symmetric convex
constraints on linear evaluations of the centered Gaussian path.
There are finitely many events for this fixed input and $T$.
Lemma~\ref{lem:gaussian-path-correlation} therefore multiplies their
probability lower bounds.  In particular, its finite-grid argument and
the continuity of the paths apply even when some constraints are
duplicates.  The strict slack just chosen ensures that their intersection
is an event of survival in the open domain $\mathcal D$.

It follows that
\begin{align}
 -\log\mathbb P\{U([0,T])\subset\mathcal D\}
 &\le C\sum_{\theta\in\Theta}
       (1+T\theta)(1+E_\theta(0,n))
       +C\sum_{j=1}^m(1+TQ_{jj})\notag\\
 &\le A(Q,R,d,n)+C'T\tr Q.\label{eq:survival-rate-bound}
\end{align}
Here \eqref{eq:scale-energy} and \eqref{eq:used-scales} bound the
coefficient of $T$ by an absolute constant times $\tr Q$, and
\[
 A(Q,R,d,n)
 =C\left[\sum_{\theta\in\Theta}(1+E_\theta(0,n))+m\right]<\infty.
\]
This additive constant may depend on the input and on the smallest
entry of $Q$; no uniform bound on it is required.

By Lemma~\ref{lem:ou-energy},
\[
 \lambda_Q^{\mathrm{OU}}(\mathcal D)
 \le C'\tr Q+\frac{A(Q,R,d,n)}{T}.
\]
Let $T$ tend to infinity while keeping the input fixed.  This removes
the additive constant, including the number of spectral scales.
The second inequality in \eqref{eq:ou-comparison} now gives
\[
 \lambda_Q(\mathcal D)\le(C'+1/2)\tr Q.
\]
Thus $H_0=C'+1/2$ is absolute.  The constants $A_0,B_0,H_0$ are
independent of $d,n,Q$ and of the recursion, as required.
\end{proof}
